% Prefer local vendored class/packages before system TeX tree.
\makeatletter
\def\input@path{{vendor/moderncv/}{assets/fonts/fontawesome7/}}
\makeatother
\documentclass[12pt,legalpaper,sans]{vendor/moderncv/moderncv} %a4paper
% SET UP BOLD NAME %
\usepackage{xstring}
\usepackage{etoolbox}
\newboolean{bold}
\newcommand{\makeauthorsbold}[1]{%
  \DeclareNameFormat{author}{%
  \setboolean{bold}{false}%
    \renewcommand{\do}[1]{\expandafter\ifstrequal\expandafter{\namepartfamily}{####1}{\setboolean{bold}{true}}{}}%
    \docsvlist{#1}%
    \ifthenelse{\value{listcount}=1}
    {%
      {\expandafter\ifthenelse{\boolean{bold}}{\mkbibbold{\namepartfamily\addcomma\addspace \namepartgiveni}}{\namepartfamily\addcomma\addspace \namepartgiveni}}%
    }{\ifnumless{\value{listcount}}{\value{liststop}}
      {\expandafter\ifthenelse{\boolean{bold}}{\mkbibbold{\addcomma\addspace \namepartfamily\addcomma\addspace \namepartgiveni}}{\addcomma\addspace \namepartfamily\addcomma\addspace \namepartgiveni}}%
      {\expandafter\ifthenelse{\boolean{bold}}{\mkbibbold{\addcomma\addspace \namepartfamily\addcomma\addspace \namepartgiveni\addcomma\isdot}}{\addcomma\addspace \namepartfamily\addcomma\addspace \namepartgiveni\addcomma\isdot}}%
      }
    \ifthenelse{\value{listcount}<\value{liststop}}
    {\addcomma\space}{}
  }
}

% CHANGE BIB BACKEND TO SORT BY YEAR
\usepackage[backend=biber, style=nature, sorting=none, maxbibnames=99,defernumbers=true]{biblatex}
%\usepackage{multibib}
\addbibresource{publications.bib}
\defbibfilter{pub_first}{keyword=published and keyword=firstauthor}
\defbibfilter{pub_co}{keyword=published and keyword=coauthor}
\defbibfilter{sub_first}{keyword=submitted and keyword=firstauthor}
\defbibfilter{sub_co}{keyword=submitted and keyword=coauthor}

% Toggle display of per-paper role annotations in bibliography.
\newif\ifshowpubroles
%\showpubrolestrue
% \showpubrolesfalse % uncomment to hide all role annotations

% Per-paper role statement for award applications.
% Add `annotation = {...}` to any @article entry in publications.bib to override
% the default role text.
\newbibmacro*{paperrolestatement}{%
  \printfield{annotation}%
}
\renewbibmacro*{finentry}{%
  \ifentrytype{article}{%
    \ifshowpubroles
      \iffieldundef{annotation}{%
        \finentry
      }{%
        \setunit{\finentrypunct\par\smallskip}%
        \printtext{\small\textit{Role: \usebibmacro{paperrolestatement}}}%
        \finentry
      }%
    \else
      \finentry
    \fi
  }{%
    \finentry
  }%
}

\moderncvstyle{casual}
\moderncvcolor{burgundy}                               % 
% Ensure section headers and photo frame always follow \moderncvcolor{...}.
\colorlet{sectioncolor}{color1}
\colorlet{subsectioncolor}{color1}
\colorlet{bodyrulecolor}{color1}
\colorlet{pictureframecolor}{color1}

% Compact the casual-style footer.
\setlength{\footskip}{70pt}
\renewcommand*{\recomputefootlengths}{\setlength{\footwidth}{\textwidth}}
\patchcmd{\makecvfoot}{\vspace{\baselineskip}}{\vspace{.2\baselineskip}}{}{}

% character encoding
%\usepackage[utf8]{inputenc}                       % if you are not using xelatex ou lualatex, replace by the encoding you are using

% adjust the page margins
\usepackage[
  left=1.06in,
  right=1.06in,
  top=1.15in,
  bottom=1.25in
]{geometry}
%\setlength{\hintscolumnwidth}{3cm}                % if you want to change the width of the column with the dates
%\setlength{\makecvheadnamewidth}{10cm}            % for the 'classic' style, if you want to force the width allocated to your name and avoid line breaks. be careful though, the length is normally calculated to avoid any overlap with your personal info; use this at your own typographical risks...

% font loading
% for luatex and xetex, do not use inputenc and fontenc
% see https://tex.stackexchange.com/a/496643
\ifxetexorluatex
  \usepackage{fontspec}
  \usepackage{unicode-math}
  \defaultfontfeatures{Ligatures=TeX}
  \setmainfont{Latin Modern Roman}
  \setsansfont{Latin Modern Sans}
  \setmonofont{Latin Modern Mono}
  \setmathfont{Latin Modern Math} 
\else
  \usepackage[utf8]{inputenc}
  \usepackage[T1]{fontenc}
  \usepackage{lmodern}
\fi

% Social icon overrides (prefer official brand icons when FA7 is fully available)
\newif\ifusefaSeven
\usefaSevenfalse
\IfFileExists{assets/fonts/fontawesome7/fontawesome7.sty}{%
  \IfFileExists{assets/fonts/fontawesome7/FontAwesome7Free-Regular-400.otf}{%
    \usefaSeventrue
  }{}%
}{}
\makeatletter
\@ifpackageloaded{fontawesome5}{\usefaSevenfalse}{}
\makeatother
\ifusefaSeven
  \usepackage{assets/fonts/fontawesome7/fontawesome7}
\fi
\definecolor{xbrand}{HTML}{000000}
\definecolor{blueskybrand}{HTML}{1185FE}
\makeatletter
\@ifundefined{faIcon}
  {\renewcommand*{\twittersocialsymbol}{{\small\textbf{X}}~}}
  {\renewcommand*{\twittersocialsymbol}{{\color{xbrand}\small\faIcon{x-twitter}}~}}
\@ifundefined{faIcon}
  {\newcommand*{\blueskysocialsymbol}{{\small\textbf{B}}~}}
  {\newcommand*{\blueskysocialsymbol}{{\color{blueskybrand}\small\faIcon{bluesky}}~}}
\makeatother

% personal data
\name{Steven Lee}{Meisler}
\title{Curriculum Vitae}                               % optional, remove / comment the line if not wanted
\address{Richards Building, 5th Floor. 3700 Hamilton Walk.}{Philadelphia, PA 19104}{USA}% optional, remove / comment the line if not wanted; the "postcode city" and "country" arguments can be omitted or provided empty
%\phone[mobile]{+1~(234)~567~890}                   % optional, remove / comment the line if not wanted; the optional "type" of the phone can be "mobile" (default), "fixed" or "fax"
\email{Steven.Meisler@pennmedicine.upenn.edu}                               % optional, remove / comment the line if not wanted
\homepage{www.stevenmeisler.com}                         % optional, remove / comment the line if not wanted
% Social icons
\social[orcid]{0000-0002-8888-1572}                  % optional, remove / comment the line if not wanted
\social[googlescholar]{tBFKYuYAAAAJ}%{0lB9-GcAAAAJ&hl}            % optional, remove / comment the line if not wanted
\social[linkedin]{Steven-Meisler}                        % optional, remove / comment the line if not wanted
\social[twitter]{StevenMeisler}                             % optional, remove / comment the line if not wanted
\social[bluesky][bsky.app/profile/stevenmeisler.com]{@StevenMeisler.com}
%\social[mastodon]{mastodon.online/@smeisler}                             % optional, remove / comment the line if not wanted
\social[github]{Smeisler}                              % optional, remove / comment the line if not wanted
%\social[gitlab]{jdoe}                              % optional, remove / comment the line if not wanted
\photo[70pt][0.4pt]{assets/images/ny_pic_head.png}
%\quote{}                                 % optional, remove / comment the line if not wanted

% bibliography adjustments (only useful if you make citations in your resume, or print a list of publications using BibTeX)
%   to show numerical labels in the bibliography (default is to show no labels)
%\makeatletter\renewcommand*{\bibliographyitemlabel}{\@biblabel{\arabic{enumiv}}}\makeatother
\renewcommand*{\bibliographyitemlabel}{[\arabic{enumiv}]}
%   to redefine the bibliography heading string ("Publications")
%\renewcommand{\refname}{Articles}

% bibliography with multiple entries
%\usepackage{multibib}
%\newcites{book,misc}{{Books},{Others}}
%----------------------------------------------------------------------------------
%            content
%----------------------------------------------------------------------------------
\begin{document}
\makeauthorsbold{Steven,Meisler}
% ---------------------------------------------------------
\makecvtitle
Steven is a developmental neuroinformaticist studying white matter maturation and its role in emerging psychopathology. As a NICHD K99/R00 postdoctoral fellow at the University of Pennsylvania, he leverages large-scale datasets and advanced analytic methods to characterize white matter developmental trajectories across infancy and adolescence. His graduate work integrated multimodal neuroimaging (fMRI, DWI) with psychocognitive assessments to investigate the neural bases of reading ability, disability, and intervention in children. Beyond his primary research, Steven is an active advocate for open and reproducible science. He moderates the NeuroStars.org support forum and teaches researchers how to process neuroimaging data using open-source, state-of-the-art tools. He received his bachelor’s and master’s degrees in bioengineering from the University of Pennsylvania (2017, 2018) and his PhD in Speech and Hearing Bioscience and Technology from Harvard University (2024). \textbf{He is on the tenure-track academic job market for positions beginning fall 2027.}
\section{Education}
\cventry{2019--2024}{PhD in Speech and Hearing Bioscience and Technology}{Harvard University, Division of Medical Sciences}{Cambridge, MA}{}{Secondary Field in Mind, Brain \& Behavior \\ Dissertation: "White Matter Structural Correlates of Reading Abilities, Disabilities, and Intervention" - Advisor: John Gabrieli, PhD}
\cventry{2017--2018}{MSE in Bioengineering}{University of Pennsylvania}{Philadelphia, PA}{\textit{3.79/4}}{Thesis: "Evaluating the Effect of Intracranial EEG Data Cleaning on Univariate and Multivariate Classifier Performance" - Advisor: Michael Kahana, PhD}
\cventry{2013--2017}{BSE in Bioengineering}{University of Pennsylvania}{Philadelphia, PA}{\textit{3.70/4}}{\textit{magna cum laude} \\ Minors in Mathematics, Jazz Studies, and Engineering Entrepreneurship}  % arguments 3 to 6 can be left empty

% ---------------------------------------------------------

\section{Research Positions}

\cventry{2024--}{Postdoctoral Fellow}{Lifespan Informatics \& Neuroimaging Center}{University of Pennsylvania (PI: Theodore Satterthwaite, MD), NIH \textbf{K99/R00} and T32 Fellow}{}{\begin{itemize}
\item Investigating white matter micro- and macrostructural trajectories in infants and toddlers, and how they are impacted by perinatal insult (e.g., low gestational age).
\item Gauging the generalizability of white matter growth curves derived from research-quality datasets to data acquired clinically.
\item Contributing to neuroimaging preprocessing software development (\textit{QSIPrep}, \textit{QSIRecon}, \textit{XCP-D}).
\end{itemize}}

\cventry{Summer 2024}{Research Scientist Intern}{Turing Medical}{St. Louis, MO}{}{
\begin{itemize}
\item Developed and tested neuroimaging software solutions for clinical practice.
\end{itemize}
}

\cventry{2019--2024}{PhD Candidate}{Gabrieli Lab}{MIT Brain and Cognitive Science (PI: John Gabrieli, PhD)}{NIH T32 Trainee and F31 Fellow}{
%Detailed achievements:
\begin{itemize}
\item Used advanced diffusion-weighted imaging models to study white matter microstructural correlates of reading in large pediatric cohorts (> 1,000 participants).
\item Studied longitudinal microstructural changes in white matter associated with response to reading intervention.
\item Examined contributions of right hemispheric activity to reading improvement after intervention among dyslexic children.
\end{itemize}
}

%\cventry{2018--2019}{Clinical Research Coordinator}{Laboratory for NeuroImaging of Coma and Consciousness}{Massachusetts General Hospital (PI: Brian Edlow, MD)}{}{
%Detailed achievements:
%\begin{itemize}
%\item Used machine learning to test if EEG and MRI biometrics in the ICU can predict long-term recovery outcomes from traumatic brain injury.
%\item Was responsible for screening patients, as well as submitting and maintaining regulatory documentation.
%\end{itemize}
%}

%\cventry{2017--2018}{Graduate Research Assistant}{Computational Memory Lab}{University of Pennsylvania (PI: Michael Kahana, PhD)}{}{
%Detailed achievements:
%\begin{itemize}
%\item Researched efficacy of various data preprocessing methods %on classification of neural activity during memory tasks
%\item Created a Python‐based GUI that visualizes and annotates EEG signals for analysis in a multi‐site DARPA‐funded clinical study
%\end{itemize}
%}
% ---------------------------------------------------------
\section{Funding}
\subsection{External}
\cventry{07/2026 -- 06/2028}{Generalizable Growth Charts of Thalamocortical Development in Early Childhood}{}{\newline Funding Agency: NIH, NICHD (type: K99/R00; HD122318)}{Role: PI}{}
\cventry{9/2025 -- 8/2026}{Hartwell Fellowship}{}
{\newline Funding Agency: Hartwell Foundation}{}{}
\cventry{8/2024 -- 8/2025}{Psychosis: A Convergent Neuroscience Perspective}{}
{\newline Funding Agency: NIH, NIMH (type: T32; MH019112)}{Role: Trainee}{(PI: Raquel E. Gur)}
\cventry{9/2023 -- 05/2024}{Neurocognitive Mechanisms of Positive Intervention Response in Reading Disability}{}{\newline Funding Agency: NIH, NICHD (type: F31; HD111139)}{Role: PI}{}
\cventry{8/2019 -- 8/2022}{Training in Speech and Hearing Sciences}{}{\newline Funding Agency: NIH, NIDCD (type: T32; DC000038)}{Role: Trainee}{(PI: Gwenaelle S. Géléoc)}
% ---------------------------------------------------------
%\newpage

% Publications from a BibTeX file without multibib
%  for numerical labels: \renewcommand{\bibliographyitemlabel}{\@biblabel{\arabic{enumiv}}}% CONSIDER MERGING WITH PREAMBLE PART
%  to redefine the heading string ("Publications"): \renewcommand{\refname}{Articles}
%\nocite{*}
%\bibliographystyle{plain}
%\bibliography{publications}

\nocite{*}
\textit{\\ * Denotes Equal Authorship}
\printbibliography[title=Published Articles (First Author), type=article, filter=pub_first, resetnumbers=true]
\printbibliography[title=Articles Under Review / Revision (First Author), type=article, filter=sub_first, resetnumbers=true]
\printbibliography[title=Published Articles (Co-Author), type=article, filter=pub_co, resetnumbers=true]
\printbibliography[title=Articles Under Review / Revision (Co-Author), type=article, filter=sub_co, resetnumbers=true]

% ---------------------------------------------------------
\printbibliography[title=Invited Talks, type=misc, resetnumbers=true] 
%\printbibliography[title=Poster Presentations, type=unpublished, resetnumbers=true] 

% Publications from a BibTeX file using the multibib package
%\section{Publications}
%\nocitebook{book1,book2}
%\bibliographystylebook{plain}
%\bibliographybook{publications}                   % 'publications' is the name of a BibTeX file
%\nocitemisc{misc1,misc2,misc3}
%\bibliographystylemisc{plain}
%\bibliographymisc{publications}                   % 'publications' is the name of a BibTeX file



% ---------------------------------------------------------
\section{Honors \& Awards}
\cvitem{2/2026-2/2028}{\textbf{Career Development Institute in Psychiatry Fellowship}, University of Pittsburgh and Stanford University}
\cvitem{9/2025-8/2026}{\textbf{Hartwell Fellowship}, supporting early-stage biomedical research with potential to benefit children, Hartwell Foundation}
\cvitem{9/2025-1/2026}{\textbf{ABCD-Repronim Teaching Fellow}, ReproNim: A Center for Reproducible Neuroimaging Computation}
\cvitem{6/2024-06/2025}{\textbf{Reproducible Neuroimaging Fellow}, ReproNim-International Neuroinformatics Coordinating Facility (INCF)}
\cvitem{5/2024}{\textbf{Ragnar \& Margaret Naess Award} for exceptional musical talent and commitment to performance, MIT}
\cvitem{8/2022}{\textbf{Patrick J. McGovern Student Travel Award}, MIT}
\cvitem{Summer 2022}{\textbf{Selected for Neurohackademy 2022}}
\cvitem{3/2022}{\textbf{Certificate of Distinction in Teaching}, Harvard University}
\cvitem{3/2020}{\textbf{Mind Brain Behavior Graduate Student Award}, Harvard University}
\cvitem{3/2019}{\textbf{University Fellowship} \textit{(declined)}, The Ohio State University}
\cvitem{2017}{\textbf{\textit{magna cum laude}}, University of Pennsylvania}
\cvitem{2015 -- 2017}{\textbf{Dean's List}, University of Pennsylvania}
% ---------------------------------------------------------

\section{Skills}
\cvitem{\textbf{Computing}}{Python, MATLAB, Bash / Terminal, R, \LaTeX}
\cvitem{\textbf{Methods}}{DWI, fMRI, EEG Time Series Analysis}
\cvitem{\textbf{Other Skills}}{EPIC, REDCap, IRB, GitHub, BIDS-apps}
% ---------------------------------------------------------
\section{Verified Peer Reviews}
\cvitem{-}{\textbf{\textit{Nature Methods, Proceedings of the National Academy of Sciences, Journal of Neuroscience, Imaging Neuroscience, Biological Psychiatry, NeuroImage, Brain Structure \& Function, npj Science of Learning, Neurobiology of Language, Brain Research Bulletin, Neuroinformatics, Neuropsychologia, Frontiers in Neuroscience}}}
%\cvitem{-}{\textbf{\textit{Journal of Neuroscience}}}
%\cvitem{-}{\textbf{\textit{Biological Psychiatry}}}
%\cvitem{-}{\textbf{\textit{NeuroImage}}}
%\cvitem{-}{\textbf{\textit{Brain Structure \& Function}}}
%\cvitem{-}{\textbf{\textit{npj Science of Learning}}}
%\cvitem{-}{\textbf{\textit{Brain Research Bulletin}}}
%\cvitem{-}{\textbf{\textit{Frontiers in Neuroscience}}}
% ---------------------------------------------------------
%\newpage
\section{Professional \& Service Affiliations}
\cvitem{2025--Present}{Trainee Committee, Flux Congress}
\cvitem{2025--Present}{Community Liaison, Open Science Special Interest Group, Organization for Human Brain Mapping}
\cvitem{09/2024}{Promoted to Moderator, NeuroStars support forum (\url{Neurostars.org})}
%\cvitem{2021--Present}{Member, Organization for Human Brain Mapping}
\cvitem{2019--Present}{Volunteer, Virtual Bedside Concerts, Mass General Brigham}
%\cvitem{2019--Present}{Member, Cognitive Neuroscience Society}
\cvitem{2019--2021}{Co-Chair, Arts \& Neuroscience Cognition Task‐Force, American Congress of Rehabilitation Medicine}
%\cvitem{2018--2019}{Member, Society for Music Perception \& Cognition}
% ---------------------------------------------------------
\newpage
\section{Teaching Experience}
\subsection{Teaching Fellow, ABCD-ReproNim Course}
\cvitem{2025}{\textbf{ABCD-ReproNim Course}}
Supported a virtual course on responsible and reproducible ABCD data analysis, including flipped-classroom instruction, data exercises, collaborative projects, and a project week focused on open-source neuroimaging workflows.

\subsection{Teaching Fellow, Harvard University}
\cvitem{Falls 2021-2023}{\textbf{How Music Works: Engineering the Acoustical World}}
Led sections and supported hands-on projects for an interdisciplinary engineering course using music and acoustics to introduce concepts in physics, mathematics, signal processing, digital audio, and design.

%\cvitem{2020}{Undergraduate Mentoring Certification, Faculty of Arts \& Science, Division of Science}
\subsection{Teaching Assistant, University of Pennsylvania}
\cvitem{Fall 2017}{\textbf{Differential Equations \& Linear Algebra}}
Supported instruction in matrices, linear systems, vector spaces, linear transformations, and systems of ordinary differential equations, with emphasis on engineering applications.

\cvitem{Fall 2016 \\ \& 2017}{\textbf{Bioengineering Signals \& Systems}}
Supported instruction in continuous and discrete signals, linear time-invariant systems, Fourier analysis, filtering, sampling, and biomedical applications including ECG, blood pressure, auditory coding, cochlear implants, and imaging.



%\section{Software}
%\subsection{\textit{FSuB-Extractor}}
%\cvitem{(Python\slash Bash)}{\url{https://github.com/smeisler/fsub_extractor}}
%This is a flexible open-source software toolbox for finding components of white matter bundles that connect to functional regions of interest. Studying these \textbf{F}unctional \textbf{Su}b-\textbf{B}undles (FSuBs) could lead to more precise studies relating brain structure, function, and behavior.
%\cvitem{2020}{Undergraduate Mentoring Certification, Faculty of Arts \& Science, Division of Science}
%\subsection{Teaching Assistant, University of Pennsylvania}
%\cvitem{Fall 2017}{Differential Equations \& Linear Algebra}
%\cvitem{Fall 2016 \\ \& 2017}{Bioengineering Signals \& Systems}

%\section{DEJI Experience}
%\cvitem{2024}{University of Rhode Island Diversity and Inclusion Badge Program}

\clearpage
\end{document}
